\chapter{Introduction}
\label{ch:introduction}

\section{Context and Motivation}
Sepsis, defined by the Sepsis-3 consensus as life-threatening organ dysfunction
caused by a dysregulated host response to infection \cite{singer2016sepsis3},
remains a major driver of mortality in the intensive care unit (ICU). Because
timely treatment materially improves outcomes, predicting the \emph{onset} of
sepsis ahead of clinical recognition is of direct clinical value. Two modelling
traditions address this problem. The first, and by far the more actively
published, is machine learning; the second, older and still foundational, is
classical statistics, namely Cox proportional-hazards regression, logistic
regression, and trajectory modelling. Statistical models retain a specific
advantage for real-time clinical use: they yield directly interpretable hazard
and odds ratios, they can be checked against formal assumptions, and they can be
benchmarked transparently against established rule-based scores such as the
Systemic Inflammatory Response Syndrome (SIRS) criteria, the quick Sequential
Organ Failure Assessment (qSOFA), and the National Early Warning Score (NEWS).

On the MIMIC-IV database \cite{johnson2023mimiciv}, however, the statistical
literature has developed unevenly. As detailed in Chapter~\ref{ch:review}, most
studies predict mortality \emph{after} sepsis is diagnosed rather than onset
itself; the proportional-hazards assumption is rarely tested before a Cox model
is fitted; and external validation is largely absent from the onset-prediction
literature specifically. The single statistical real-time onset benchmark,
TREWScore \cite{henry2015trewscore}, was fitted on the older MIMIC-II database
in 2015 and has not been revisited on MIMIC-IV with purely statistical,
interpretable tools.

\section{Research Gap}
\label{sec:gap}
Three gaps motivate this work and are established in detail in
Chapter~\ref{ch:review}:
\begin{itemize}
  \item \textbf{G1 (Onset versus prognosis).} MIMIC-IV statistical work
  overwhelmingly predicts mortality after diagnosis. Genuine pre-onset,
  real-time detection with purely statistical tools has not been revisited on
  MIMIC-IV since TREWScore's 2015 MIMIC-II study.
  \item \textbf{G2 (Untested proportional hazards).} Sepsis hazards are
  frequently non-proportional \cite{kasal2004cox}, yet no reviewed MIMIC-IV Cox
  study reports testing the assumption, weakening confidence in the reported
  hazard ratios.
  \item \textbf{G3 (Validation asymmetry).} External validation is common in the
  prognostic nomogram literature but essentially absent from onset-prediction
  and trajectory studies, precisely where systematic reviews show performance
  degrades most \cite{wang2025srpm}.
\end{itemize}

\section{Research Question and Objectives}
\label{sec:rq}
These gaps motivate one primary research question and three sub-questions.

\vspace{4pt}
\noindent\textbf{RQ.} \emph{Can a purely statistical (Cox proportional-hazards /
time-varying / landmark) model, using routinely available real-time vitals and
labs in MIMIC-IV, predict impending Sepsis-3 onset ahead of clinical
recognition, with performance comparable to or exceeding the rule-based scores
SIRS, qSOFA, and NEWS?}

\begin{itemize}
  \item \textbf{SQ1.} Which routinely measured variables are the strongest
  statistically significant predictors of onset, and does the
  proportional-hazards assumption hold for the fitted Cox model?
  \item \textbf{SQ2.} Does modelling trajectories through a time-varying or
  landmark extension improve early detection over a static single-time-point
  model?
  \item \textbf{SQ3.} How does discriminative performance compare with SIRS,
  qSOFA, and NEWS, and does it survive the transition from internal to external
  (eICU-CRD) evaluation?
\end{itemize}

\section{Contributions}
\label{sec:contrib}
This thesis makes the following contributions, each mapped to the gap it
closes.
\begin{enumerate}
  \item \textbf{C1: First interpretable statistical real-time Sepsis-3 onset
  model on MIMIC-IV (closes G1).} Continuously updated hourly risk is reached
  through a dynamic landmark supermodel
  \cite{vanhouwelingen2007landmark,fries2023dynamiclm}, and incident sepsis is
  explicitly separated from prevalent sepsis using an hour-six landmark.
  \item \textbf{C2: Explicit proportional-hazards testing with a triggered
  time-varying reformulation (closes G2).} The Schoenfeld residual test is
  reported for the onset Cox model; because the global test fails, a
  time-varying Cox model is fitted in direct response.
  \item \textbf{C3: Reframing the task.} On the correct incident-onset
  target, SOFA is inversely discriminative, qSOFA is near chance, and SIRS is
  weak, while a multivariable landmark model reaches an area under the curve of
  0.775, showing that severity scores detect established organ failure, not
  impending onset.
  \item \textbf{C4: Informative-missingness characterisation.} Binary
  ``was this lab ordered'' indicators are stronger predictors than the imputed
  values themselves \cite{groenwold2020missingness,singh2021missingness},
  quantifying a label-recognition ceiling.
  \item \textbf{C5: A reproducible, transportable pipeline that is externally
  validated on eICU-CRD (closes G3).} Because the models are closed-form, their
  coefficients transport directly to a second database; the frozen model is
  applied without refitting to the eICU Collaborative Research Database
  \cite{pollard2018eicu}.
\end{enumerate}

\section{Thesis Structure}
Chapter~\ref{ch:methodology} sets out the research and software methodology.
Chapter~\ref{ch:review} reviews the related statistical and machine-learning
literature and formalises the gaps above. Chapter~\ref{ch:design} describes the
system architecture and the design of each pipeline stage.
Chapter~\ref{ch:evaluation} presents and critically discusses the internal and
external results. Chapter~\ref{ch:conclusion} concludes and identifies avenues
for further work.
